\emph{Preliminaries.} By Theorems \claimref{synthesis.signed-directions-and-the-integer-clifford-algebra} and \claimref{synthesis.division-gives-the-rationals} the blades are a basis over $k$, $\dagger$ (inversion on the group of signed directions, extended linearly) is an anti-automorphism, and $\langle x^{\dagger}x\rangle_{0}$ is the sum of the squares of the coefficients of $x$, so states exist ($x \mapsto \langle x\rangle_{0}$ is one). For a state and self-adjoint $h$, $0 \le \omega((h - t)(h - t)^{\dagger}) = \omega(h^{2}) - 2t\omega(h) + t^{2}$ at $t = \omega(h)$ gives $\omega(h)^{2} \le \omega(h^{2})$; and $\omega(h^{2}) = \omega(hh^{\dagger}) \ge 0$.

\emph{(i).} $A_{+}A_{-} = A_{1}^{2} - A_{1}A_{2} + A_{2}A_{1} - A_{2}^{2} = -[A_{1},A_{2}]$ and $A_{-}A_{+} = [A_{1},A_{2}]$; $A_{\pm}^{2} = 2 \pm (A_{1}A_{2} + A_{2}A_{1})$. Regrouping, $C_{AB} = (A_{1} + A_{2})B_{1} + (A_{1} - A_{2})B_{2}$, and likewise for $D$. Put $U = A_{+}B_{1}$, $V = A_{-}D_{2}$. Since $B_{1}$ and $D_{2}$ commute with $A_{1}, A_{2}$ and with each other, $U^{\dagger} = B_{1}A_{+} = U$, $V^{\dagger} = V$, $U^{2} = A_{+}^{2}$, $V^{2} = A_{-}^{2}$, and $UV + VU = (A_{+}A_{-} + A_{-}A_{+})B_{1}D_{2} = 0$. So $(\lambda U + \mu V)^{2} = \lambda^{2}A_{+}^{2} + \mu^{2}A_{-}^{2}$. The same holds for $U' = A_{-}B_{2}$, $V' = A_{+}D_{1}$ with $A_{+}$ and $A_{-}$ exchanged.

\emph{(ii).} Let $a = \omega(A_{+}^{2})$; then $a \ge 0$ and $4 - a = \omega(A_{-}^{2}) \ge 0$. Write $u = \omega(U)$, $v = \omega(V)$, $u' = \omega(U')$, $v' = \omega(V')$, so that $\omega(C_{AB}) = u + u'$ and $\omega(C_{AD}) = v' + v$. By (i) and the preliminaries, for all scalars $\lambda, \mu$,
\[
(\lambda u + \mu v)^{2} \le \lambda^{2}a + \mu^{2}(4 - a), \qquad (\lambda u' + \mu v')^{2} \le \lambda^{2}(4 - a) + \mu^{2}a . \tag{$*$}
\]
\emph{Case $0 < a < 4$.} Put $\tau = u^{2}/a + v^{2}/(4 - a)$ and take $\lambda = u/a$, $\mu = v/(4 - a)$ in the first inequality: $\tau^{2} \le \tau$, so $\tau \le 1$. Likewise $\tau' = u'^{2}/(4 - a) + v'^{2}/a \le 1$. The identity
\[
4\Bigl(\frac{u^{2}}{a} + \frac{u'^{2}}{4 - a}\Bigr) - (u + u')^{2} = \frac{\bigl((4 - a)u - au'\bigr)^{2}}{a(4 - a)} \ge 0
\]
and the same identity for $v'$, $v$ (with $a$ under $v'$) give $(u + u')^{2} + (v + v')^{2} \le 4(\tau + \tau') \le 8$. \emph{Case $a = 0$.} From ($*$): $\mu = 0$ gives $u = 0$, and $\lambda = 0$ gives $v^{2} \le 4$; in the second, $\lambda = 0$ gives $v' = 0$ and $\mu = 0$ gives $u'^{2} \le 4$. So the sum of squares is $u'^{2} + v^{2} \le 8$. \emph{Case $a = 4$} is the same with the roles exchanged. If both values exceeded $2$ in absolute value the sum of squares would exceed $8$. For a single pair: ($*$) with $\mu = 0$ gives $u^{2} \le a$ and $u'^{2} \le 4 - a$, and the identity above gives $(u + u')^{2} \le 8$ when $0 < a < 4$, while for $a \in \{0, 4\}$ one of $u, u'$ vanishes and the other has square at most $4$.

\emph{(iii).} If $A_{1}$ and $A_{2}$ commute then $A_{+}A_{-} = 0 = A_{-}A_{+}$, so $C_{AB}^{2} = A_{+}^{2}B_{1}^{2} + A_{-}^{2}B_{2}^{2} + A_{+}A_{-}B_{1}B_{2} + A_{-}A_{+}B_{2}B_{1} = A_{+}^{2} + A_{-}^{2} = 4$, and $\omega(C_{AB})^{2} \le \omega(C_{AB}^{2}) = 4$. In the commutative case each observable is $\sum_{i}\epsilon_{i}q_{i}$ with $\epsilon_{i}^{2} = 1$ (its square is $1 = \sum q_{i}$), so $C_{AB} = \sum_{i}\gamma_{i}q_{i}$ with $\gamma_{i} = \alpha_{1}(\beta_{1} + \beta_{2}) + \alpha_{2}(\beta_{1} - \beta_{2})$ for signs $\alpha, \beta$; one of $\beta_{1} \pm \beta_{2}$ is $0$ and the other $\pm 2$, so $\gamma_{i} = \pm 2$. $\omega(q_{i}) = \omega(q_{i}q_{i}^{\dagger}) \ge 0$ and $\sum_{i}\omega(q_{i}) = 1$. For commuting observables $A$, $B$ the product $AB$ is an observable, so $\omega(AB)^{2} \le \omega(1) = 1$, and $|\omega(C_{AB})| \le 4$ by the triangle inequality; (ii) gives $\omega(C_{AB})^{2} \le 8 < 16$.

\emph{(iv).} The adjoint fixes $e_{j}$ and negates $f_{j}$; $p_{j}$ is self-adjoint of square $1$ and anticommutes with $e_{j}$, $f_{j}$, and commutes with the generators of the other modes. By clause (ii) of Theorem \claimref{synthesis.mode-split-of-the-clifford-algebra-and-its-klein-twist-to-the-tensor-product} with untwisted set $\{1\}$ in the algebra of modes $1, 2$, the element $e_{2}p_{1}$ has square $1$ and commutes with the algebra of mode $1$; with untwisted set $\{1, 2\}$ in $\Cl(3,3)$, $e_{3}p_{1}p_{2}$ has square $1$ and commutes with the algebra of modes $1, 2$, which contains $e_{1}, p_{1}, e_{2}p_{1}, p_{2}$; $p_{3}$ is even in mode $3$ and commutes with modes $1, 2$. Both twisted elements are self-adjoint (a product of commuting self-adjoint factors). $e_{3}p_{1}$ anticommutes with $e_{2}p_{1}$, since $e_{3}$ and $e_{2}$ anticommute. The total parity $p_{1}p_{2}p_{3}$ anticommutes with $e_{1}$, $e_{2}p_{1}$ and $e_{3}p_{1}p_{2}$. The three factors of $\rho$ are $1 + g$ for commuting observables $g$ (the check verifies this and the values), so $\rho$ is a positive multiple of a self-adjoint idempotent; with $B_{1} = B_{2}$, $C_{AB} = 2A_{1}B_{1}$ and $\omega_{\rho}(A_{1}B_{1}) = 1$, likewise for $D$. The last statement is (ii) with $S = \tfrac{14}{5}$; the value $\tfrac{14}{5}$ and the sample of directions are computed by the check.

\emph{(v).} The first statement is (ii). The correlations of the pair state are computed from its expression in clause (v) of the mode-split Theorem; that expression is $\rho = \rho^{(0)} + r\rho^{(1)} + cs\,\rho^{(2)}$ with fixed elements $\rho^{(i)}$, so each correlation is $\alpha + \beta r + \gamma\,cs$ with rational constants, and the check evaluates all four at three weight pairs whose vectors $(1, r, cs)$ are linearly independent, which determines the constants: $2cs$, $1$, $0$, $0$ (these are the values $2cs$ and $1$ are those of the amplitudes $c$ on $\ket{00}$ and $s$ on $\ket{11}$: $e_{1}\,e_{2}p_{1}$ exchanges the two terms and $p_{1}p_{2}$ fixes both). With $B_{\pm} = zp_{2} \pm x\,e_{2}p_{1}$: $C = p_{1}(B_{+} + B_{-}) + e_{1}(B_{+} - B_{-}) = 2z\,p_{1}p_{2} + 2x\,e_{1}\,e_{2}p_{1}$, of value $2z + 4cs\,x$; by the Cauchy inequality $(z + 2cs\,x)^{2} \le (z^{2} + x^{2})(1 + 4c^{2}s^{2})$, and $4c^{2}s^{2} = 1 - r^{2}$. The numbers at the stated point are arithmetic.
